\section{Why PAC-Level Examples Exist Anyway}
\label{sec:why-pac-examples}

An example that pokes registers directly is an antipattern in a project that has a HAL. It duplicates
knowledge, it names pins inline, and it is exactly the code a reader should not copy into a product.
This project keeps nineteen of them, deliberately, and this section is the argument.

\subsection{The Ladder}

The same demo exists at four levels, and the differences are the whole curriculum:

\begin{table}[htbp]
    \centering
    \footnotesize
    \begin{tabular}{@{}p{0.26\textwidth}p{0.12\textwidth}p{0.50\textwidth}@{}}
        \toprule
        \textbf{Example} & \textbf{Lines} & \textbf{What a reader has to know to change it} \\
        \midrule
        \codeinline{pac\_led\_blink}   & ~90 & the port's MODEL and DOUT registers, and its clock gate \\
        \codeinline{hal\_led\_blink}   & ~50 & that a pin has a port, a bit and a polarity \\
        \codeinline{module\_led\_blink}& ~35 & that the board has an LED \\
        \codeinline{base\_led\_blink}  & ~35 & that the module is on a carrier \\
        \bottomrule
    \end{tabular}
    \caption{One program, four rungs. Each rung removes a fact the reader must hold in their head —
    and each rung is only trustworthy if the rung below it is still testable.}
    \label{tab:blink-ladder}
\end{table}

\subsection{What the PAC Level Proves}

The bottom rung is not a tutorial. It is an instrument:

\begin{itemize}
    \item \textbf{It proves the PAC matches the silicon.} If \codeinline{pac\_led\_blink} blinks, the
          SVD, the generator and the part agree. If it does not, no amount of HAL work will help,
          and the failure has exactly one suspect.
    \item \textbf{It localises a regression.} When a HAL example stops working, the PAC example says
          whether the hardware path or the abstraction broke. That question has cost this project a
          full session at least once.
    \item \textbf{It documents facts the HAL deliberately hides.} \codeinline{pac\_pullup\_probe}
          reads \codeinline{PORTB\_DIN} with every pin in input mode and proves the external pull-ups
          exist; \codeinline{pac\_i2c\_scan} proves the on-chip controller cannot complete a read.
          Neither is expressible through the HAL, because the HAL exists to make those facts
          invisible.
\end{itemize}

\subsection{The Rule That Keeps It Honest}

Proof-of-concept code rots faster than any other code, so it is held to a stricter rule than the
HAL: \textbf{a PAC example may only name the layer it demonstrates}.

\begin{minted}{bash}
# from this project's gate: no example outside the PAC crate may mention the PAC
grep -rln 'efr32mg24_pac\|pac::' crates/*/examples/ \
  | grep -v '^crates/efr32mg24-pac/'
\end{minted}

The inverse is enforced too: a PAC example that imports the HAL is not a PAC example, and would
quietly stop testing what it claims to test. Between the two checks, each example belongs to exactly
one rung of the ladder, and a reader can trust which rung they are standing on.

\subsection{When to Delete One}

\begin{itemize}
    \item When its fact is now proved elsewhere \emph{and} the example no longer adds a distinct
          failure signal. Two examples that fail for the same reason are one example.
    \item Never merely because it looks old. This project retired six HAL examples that had rotted
          into non-compiling code, and kept every PAC example that still builds and still tells you
          something the layer above cannot.
\end{itemize}
